\documentclass[11pt]{article}
\usepackage[margin=0.9in]{geometry}
\usepackage{amsmath,amssymb,amsthm,booktabs}
\usepackage[T1]{fontenc}
\usepackage{lmodern,microtype}
\usepackage[colorlinks=true,linkcolor=blue,urlcolor=blue]{hyperref}
\newtheorem{theorem}{Theorem}[section]
\newtheorem{lemma}[theorem]{Lemma}
\newtheorem{proposition}[theorem]{Proposition}
\newtheorem{corollary}[theorem]{Corollary}
\theoremstyle{remark}
\newtheorem{remark}[theorem]{Remark}
\newcommand{\R}{\mathbb R}
\newcommand{\C}{\mathbb C}
\newcommand{\norm}[1]{\left\|#1\right\|}
\newcommand{\diag}{\operatorname{diag}}
\newcommand{\ran}{\operatorname{ran}}
\newcommand{\HH}{\mathcal H}
\title{Toward a four-dimensional Yang--Mills construction\\
\large Notebook 09: the flat fine inverse and a current neighborhood}
\author{Research working document}
\date{16 October 2026}
\begin{document}
\maketitle

\begin{abstract}
We prove exponential locality of the full homogeneous flat fine
inverse in the scalar gauge used by the deposited source account,
including every placement of at most two fine differences. An explicit
rank-one correction cancels its apparent massless poles. A separate
constrained scalar inverse controls the differentiated gauge vertices.
Together these estimates yield a complex neighborhood of bounded
auxiliary currents with radius independent of fine refinement and
periodic volume. The coarse precision has an exact quadratic remainder
at zero current, factorial bounds at every order, and exponential
paths through localized current supports. All operator norms use
physical cell volume. Notebook 11 separately treats the nonlinear
constraint-current term at flat background. Nonflat backgrounds and
the interacting continuum limit remain open.
\end{abstract}

\noindent\textbf{Status.}
Research proofs checked by separate internal derivations and audits.
No external certification or novelty claim is made. Operator matching
refers to the pinned deposited formulas; it does not reproduce every
unavailable primary-paper argument. Throughout, complex adjoints mean
the analytic continuation of the real bilinear transpose, rather than
conjugation of a varying complex parameter.

\section{Scope and normalization}

The dimension is fixed, with $d=4$ the intended application. Work on
the infinite homogeneous lattice or a compatible periodic regulator.
The physical coarse spacing is one and the fine spacing is $M^{-1}$,
where $M$ is odd. Fine inner products use measure $M^{-d}$ and coarse
inner products use counting measure. The normalized averages are
\[
 (Q_0\lambda)(n)=M^{-d}\sum_{s\in I_M^d}\lambda(n+s/M),
 \quad I_M=\{-(M-1)/2,\ldots,(M-1)/2\},
\]
\[
 (Q_M A)_\mu(n)=M^{-d-1}\sum_{s\in I_M^d}\sum_{j=0}^{M-1}
 A_\mu(n+(s+je_\mu)/M).
\]
Write $Q=Q_M$. The physical Maxwell precision $q_M$ is defined by
\[
 \langle B,q_MB\rangle=\inf_{QA=B}\norm{d_\eta A}_\eta^2.
\]
Notebook 07 proves its uniform holomorphic symbol and cell locality.
Notebook 08 constructs its unique source-gauge extension $H_{\rm src}$
and verifies the deposited averaging and gauge conventions. The
following proof controls the full fine inverse required to perturb
that extension. Its positive coarse penalty is an auxiliary constraint
penalty: the resulting fine inverse gap is not a physical mass gap.

We use the alias estimates and physical Fourier normalization proved
in Notebook 07. Relevant multiplier definitions and cancellations are
recalled below. Every refinement-uniform bound is for a fixed compact
positive interval of penalties, or for the source value $a=1$.
% Full flat fine inverse: fragment for Notebook 09, no standalone preamble.
% Requires amsmath/amssymb/amsthm and the existing \norm, \diag macros.
\section{The full flat source inverse}
\label{inv-section}

We now control the fine inverse itself, rather than only the minimized
coarse action or its retained-coordinate covariance. Fine spacing is
$\eta=M^{-1}$, where $M$ is odd, and the coarse physical spacing is one.
All fine pairings use volume weight $\eta^d$, while coarse pairings use
unit counting measure. Let $D$ be the fine scalar gradient,
$\mathcal L=D^*D$ its positive scalar Laplacian, and $K$ the fine
Maxwell precision. Thus $KD=0$ and $K+DD^*=\mathcal L I$.
Let $Q=Q_M^{\rm st}$ be the normalized one-form straight-string average
and $Q_0$ the centered scalar average, so that $QD=d_cQ_0$.

The scalar source gauge is
\begin{equation}\label{inv-source-projection}
 G'=(\mathcal L+Q_0^*a'Q_0)^{-1},\qquad
 \mathcal P=G'Q_0^*(Q_0G'^2Q_0^*)^{-1}Q_0G',\qquad R=I-\mathcal P.
\end{equation}
The scalar projection $\mathcal P$ must not be confused with a Fourier
denominator $P_r$ below. The positive scalar penalty $a'$ in this
definition is distinct from the one-form penalty $a$ in
\begin{equation}\label{inv-fine-precision}
 \mathcal A_a=K+DRD^*+Q^*aQ,\qquad G_0=\mathcal A_a^{-1}.
\end{equation}
We first allow a fixed positive $a$, uniformly in a fixed interval
$0<a_-\le a\le a_+<\infty$. The source convention of Raw I,
lines 13583--13593, is $a=\ell^{-2}$; in the present homogeneous
unit-scale normalization it is exactly $a=1$.

\begin{theorem}\label{inv-theorem}
The inverse in \eqref{inv-fine-precision} exists on the homogeneous
infinite lattice and every compatible periodic regulator. Its
physical-cell kernel satisfies, for fixed positive $C,\rho$,
\begin{equation}\label{inv-kernel-bound}
 \norm{(\mathcal D_\alpha G_0\mathcal D_\beta)(n,m)}
 _{\mathrm{cell},2\to\mathrm{cell},2}
 \le Ce^{-\rho|n-m|_1},\qquad \alpha+\beta\le2.
\end{equation}
Here $n,m$ are unit coarse cells, each $\mathcal D$ is a specified
composition of fine forward or backward differences of its indicated
order, and fixed curl or divergence combinations are included.
Constants are independent of odd $M$, volume, and $a\in[a_-,a_+]$.
The cell Hilbert norm is the fine physical-volume norm, including the
fixed orientation and color fibers. At every smaller positive exponent,
\eqref{inv-kernel-bound} gives a uniform exponentially weighted
cell-Schur bound. The periodic kernels are exact periodizations and
obey the same bounds in periodic distance.

If $q_M$ denotes the physical Maxwell minimum, then the penalized
coarse inverse and normalized extension satisfy
\begin{equation}\label{inv-coarse-result}
 q_0:=(QG_0Q^*)^{-1}=q_M+aI,
 \qquad G_0Q^*q_0=H_{\rm src},
\end{equation}
where $H_{\rm src}$ is the exact Maxwell minimizing extension in the
matched scalar source gauge $RD^*H_{\rm src}=0$.
\end{theorem}

The theorem is specific to the homogeneous flat source definitions.
It is not a theorem for arbitrary nonflat backgrounds, Dirichlet
domains, or mixed retained histories. Fine derivative bounds in
\eqref{inv-kernel-bound} are cell $L^2$ operator bounds, not
pointwise regularity assertions.

\subsection{Exact aliases and the penalized Feynman inverse}

Let $I_M=\{-(M-1)/2,\ldots,(M-1)/2\}$,
$p_r=z+2\pi r$, and $r\in I_M^d$. Use
\[
 u_{r,\mu}^{\pm}=M(e^{\pm ip_{r,\mu}/M}-1),\quad
 P_r=\sum_\mu u_{r,\mu}^+u_{r,\mu}^-,\quad
 d_\mu^\pm=e^{\pm iz_\mu}-1,\quad b_\mu=d_\mu^+d_\mu^-.
\]
More explicitly, the averaging multipliers are
\[
 a_r=\prod_{\nu=1}^d\left(M^{-1}\sum_{s\in I_M}
       e^{ip_{r,\nu}s/M}\right),\qquad
 \omega_{r,\mu}=a_r M^{-1}\sum_{j=0}^{M-1}e^{ip_{r,\mu}j/M}.
\]
Write $a_r=a_M(p_r)$ for the centered scalar averaging multiplier
and $\omega_{r,\mu}$ for the one-form multiplier, including its
additional longitudinal string factor. Superscript minus means
evaluation at the negative full fine momentum; it agrees with
complex conjugation on real momenta only. String telescoping gives
\begin{equation}\label{inv-cochain}
 \omega_{r,\mu}u_{r,\mu}^+=a_r d_\mu^+,
 \qquad \omega_{r,\mu}^-u_{r,\mu}^-=a_r^-d_\mu^-.
\end{equation}
Alias space has counting Hilbert norm. With physical-volume Parseval
normalization, no additional power of $M$ occurs in an alias covariance
sum.

Work initially in a real nonzero coarse fiber. The scalar source
projection is the orthogonal projection onto
$\operatorname{ran}(\mathcal L^{-1}Q_0^*)$. Therefore
\begin{equation}\label{inv-projection-alias}
 \mathcal P=\mathcal L^{-1}Q_0^*W_2^{-1}Q_0\mathcal L^{-1},
 \qquad W_2=\sum_r\frac{a_ra_r^-}{P_r^2}.
\end{equation}
This agrees with \eqref{inv-source-projection}: multiplying by the
positive scalar penalty changes the coarse parametrization of
$\operatorname{ran}(G'Q_0^*)$, but not that range. Its constant scalar
mode will be treated through the original penalized definition.

Define
\[
 G_F=\diag_r(P_r^{-1}I),\qquad
 W=QG_FQ^*=\diag_\mu(c_\mu),\qquad
 c_\mu=g_\mu/P_0+R_\mu,
\]
\[
 g_\mu=\omega_{0,\mu}\omega_{0,\mu}^-,\qquad
 R_\mu=\sum_{r\ne0}\frac{\omega_{r,\mu}\omega_{r,\mu}^-}{P_r}.
\]
The Feynman operator with the one-form penalty is
$\mathcal B=\mathcal L I+Q^*aQ$. Its inverse obeys
\begin{equation}\label{inv-F-woodbury}
 F:=\mathcal B^{-1}
 =G_F-G_FQ^*(a^{-1}I+W)^{-1}QG_F.
\end{equation}
Put
\begin{equation}\label{inv-Da}
 \widehat R_\mu=R_\mu+a^{-1},\qquad
 \mathcal D_\mu=g_\mu+P_0\widehat R_\mu,\qquad
 k_\mu=\mathcal D_\mu^{-1}.
\end{equation}
Then $(a^{-1}I+W)^{-1}=P_0\diag(k_\mu)$ and
\begin{equation}\label{inv-F-alias}
 F_{r\mu,s\nu}=\delta_{\mu\nu}
 \left[\frac{\delta_{rs}}{P_r}
 -\frac{P_0\omega_{r,\mu}^-\omega_{s,\mu}}
              {P_rP_s\mathcal D_\mu}\right].
\end{equation}
The central blocks cancel exactly to
\begin{align}
 F_{0\mu,0\nu}
 &=\delta_{\mu\nu}\widehat R_\mu/\mathcal D_\mu,
 \label{inv-F-central}\\
 F_{0\mu,s\nu}
 &=-\delta_{\mu\nu}
   \frac{\omega_{0,\mu}^-\omega_{s,\mu}}{P_s\mathcal D_\mu}
       &&(s\ne0),\nonumber\\
 F_{r\mu,0\nu}
 &=-\delta_{\mu\nu}
   \frac{\omega_{r,\mu}^-\omega_{0,\mu}}{P_r\mathcal D_\mu}
       &&(r\ne0).\nonumber
\end{align}
Use \eqref{inv-F-alias} only with both aliases high, and
\eqref{inv-F-central} otherwise. These expressions contain no
central massless denominator.

\subsection{The scalar projection and the second Woodbury identity}

The Hodge identity and $R=I-\mathcal P$ give
\begin{equation}\label{inv-rank-subtraction}
 \mathcal A_a=\mathcal B-vW_2^{-1}v^-,\qquad
 v_{r,\mu}=u_{r,\mu}^+a_r^-/P_r,
 \quad v_{r,\mu}^-=a_ru_{r,\mu}^-/P_r.
\end{equation}
Here $v^-$ is a row, the analytic transpose at negative arguments.
On real fibers it is the Hilbert adjoint of $v$.

The real fine precision is positive definite. Indeed, a vector in its
kernel has zero curl, zero coarse average, and zero source-gauge term.
The zero-curl field is $D\lambda+c$, with $c$ a constant harmonic
one-form. Taking the coarse mean of $Q(D\lambda+c)=0$ gives $c=0$;
subtracting a scalar constant makes $Q_0\lambda=0$. Then
$\mathcal L\lambda\in\operatorname{ran}R$, while the source-gauge
condition says $R\mathcal L\lambda=0$. Thus $D\lambda=0$.
This proves finite-regulator positivity with the harmonic sector
retained. The subtractive Woodbury identity therefore gives
\begin{equation}\label{inv-G-woodbury}
 G_0=F+Fv\,(W_2-v^-Fv)^{-1}v^-F.
\end{equation}

Equation \eqref{inv-cochain} implies the three exact contractions
\begin{equation}\label{inv-three-contractions}
 v^-G_Fv=W_2,\qquad QG_Fv=W_2d^+,
 \qquad v^-G_FQ^*=W_2d^-.
\end{equation}
For example, the $\mu$ component of the middle term is
$\sum_r\omega_{r,\mu}u_{r,\mu}^+a_r^-/P_r^2=W_2d_\mu^+$.
Let $\mathfrak a=a_0a_0^-$. The central cochain identity is
$\sum_\mu b_\mu/g_\mu=P_0/\mathfrak a$. Define
\begin{equation}\label{inv-Ea}
 \mathcal E_a=\mathfrak a^{-1}
 -\sum_\mu\frac{b_\mu\widehat R_\mu}{g_\mu\mathcal D_\mu}.
\end{equation}
Subtracting $1/\mathcal D_\mu-1/g_\mu$ proves
\begin{equation}\label{inv-Ea-identity}
 \sum_\mu b_\mu k_\mu=P_0\mathcal E_a.
\end{equation}
It follows that
\begin{align}
 W_2-v^-Fv
 &=W_2^2d^-(a^{-1}I+W)^{-1}d^+
  =W_2^2P_0^2\mathcal E_a.
 \label{inv-woodbury-denominator}
\end{align}

Introduce the scalar second-order remainder and its canceled denominator
\begin{equation}\label{inv-Ds2}
 R_{s,2}=\sum_{r\ne0}\frac{a_ra_r^-}{P_r^2},\qquad
 D_{s,2}=\mathfrak a+P_0^2R_{s,2},\qquad
 W_2=D_{s,2}/P_0^2.
\end{equation}
Thus \eqref{inv-woodbury-denominator} is
$D_{s,2}^2\mathcal E_a/P_0^2$. In particular it contains $W_2^2$,
not $W_2$.

\subsection{The final canceled inverse}

Set $\beta=(P_0/D_{s,2})Fv$. For a high alias,
\begin{equation}\label{inv-beta-high}
 \beta_{r,\mu}
 =\frac{P_0u_{r,\mu}^+a_r^-}{P_r^2D_{s,2}}
 -\frac{\omega_{r,\mu}^-k_\mu d_\mu^+}{P_r},
 \qquad r\ne0.
\end{equation}
For the central entry use
\[
 \omega_{0,\mu}^-k_\mu d_\mu^+
 =(u_{0,\mu}^+/a_0)g_\mu k_\mu,
\]
\[
 \mathfrak a/D_{s,2}=1-P_0^2R_{s,2}/D_{s,2},\qquad
 g_\mu k_\mu=1-P_0\widehat R_\mu k_\mu.
\]
The resulting exact central formula is
\begin{equation}\label{inv-beta-central}
 \beta_{0,\mu}=
 \frac{u_{0,\mu}^+}{a_0}
 \left(\widehat R_\mu k_\mu
                  -\frac{P_0R_{s,2}}{D_{s,2}}\right).
\end{equation}
The row $\beta^-=(P_0/D_{s,2})v^-F$ is its analytic transpose:
in the high formula it has entries
$P_0a_ru_{r,\mu}^-/(P_r^2D_{s,2})
-d_\mu^-k_\mu\omega_{r,\mu}/P_r$; centrally use
$u_{0,\mu}^-/a_0^-$ times the same bracket.
Equations \eqref{inv-G-woodbury} and
\eqref{inv-woodbury-denominator} now become
\begin{equation}\label{inv-final-formula}
 \boxed{\displaystyle G_0=F+\beta\mathcal E_a^{-1}\beta^-.}
\end{equation}
Every central pole has canceled. The update has a positive sign, and
the second term in the bracket of \eqref{inv-beta-central} has a
negative sign. The formula removes the entire complex cone $P_0=0$.

\subsection{A strip and constants uniform in refinement}

Choose $\pi<K_0<2\pi$ and set
\[
 c_* =\frac{\sin(K_0/2)}{K_0/2},\qquad
 H_* =\frac{4(2\pi-K_0)^2}{\pi^2}.
\]
On the enlarged real cube $|x_\mu|\le K_0$, the alias estimates of
Notebook 07 give
\[
 \mathfrak a\ge c_*^{2d},\quad
 g_\mu\ge g_*:=c_*^{2d+2},\quad
 0\le R_\mu\le H_*^{-1},\quad
 0\le R_{s,2}\le H_*^{-2},\quad P_0\le dK_0^2,
\]
with $g_\mu,\mathfrak a\le1$. Consequently
\[
 g_*\le\mathcal D_\mu\le
 D_*:=1+dK_0^2(H_*^{-1}+a_-^{-1}),\qquad
 c_*^{2d}\le D_{s,2}\le1+d^2K_0^4H_*^{-2}.
\]
For $x\ne0$, \eqref{inv-Ea-identity} gives
\begin{equation}\label{inv-Ea-lower}
 \mathcal E_a(x)=
 \frac{\sum_\mu b_\mu(x)/\mathcal D_\mu(x)}{P_0(x)}
 \ge c_*^2/D_*.
\end{equation}
Here $\sum b_\mu\ge c_*^2|x|^2$ and $P_0\le|x|^2$.
At zero the defining expression gives $\mathcal E_a(0)=1$.

On an initial fixed complex neighborhood, high denominators obey
$|P_r(z)|\ge c|r|^2$ for $r\ne0$, including the edge aliases.
Finite alias Parseval bounds the sums of absolute paired scalar and
one-form averaging factors independently of $M$. Hence $R_\mu$,
$R_{s,2}$, and the central multipliers are uniformly holomorphic and
bounded. The added $a^{-1}$ is bounded by the hypothesis on $a$.

Use nested real cubes with positive margins. Uniform Cauchy estimates
first allow an imaginary shrinking that keeps $\mathfrak a,g_\mu,
\mathcal D_\mu,D_{s,2}$ separated from zero. The expression
\eqref{inv-Ea} is then holomorphic and uniformly bounded. A second
shrinking, using \eqref{inv-Ea-lower}, keeps $\mathcal E_a$ separated
from zero. This gives one positive strip width independent of $M,a$,
with the real fundamental cube in the interior of each enlarged chart.
The inverse of $a_0$ is also bounded because $a_0a_0^-=\mathfrak a$
is separated from zero and both factors are bounded. Real positivity
has thus been combined with uniform analytic estimates; it is not
being used alone to claim complex invertibility.

\subsection{Zero mode and the independent $M=1$ check}

At coarse momentum zero, $a_0=\omega_{0,\mu}=1$, all high
$a_r,\omega_{r,\mu}$ vanish, and
\[
 R_\mu=R_{s,2}=0,\qquad
 \mathfrak a=g_\mu=\mathcal D_\mu=D_{s,2}=\mathcal E_a=1,
 \qquad\beta=0.
\]
The canceled formula gives
\begin{equation}\label{inv-zero-mode}
 (G_0)_{0\mu,0\nu}=a^{-1}\delta_{\mu\nu},\qquad
 (G_0)_{r\mu,s\nu}=\delta_{rs}\delta_{\mu\nu}/P_r
       \quad(r,s\ne0),
\end{equation}
with zero mixed blocks. The source projection $R$ is zero on the
constant scalar mode and identity on high scalar aliases. Hence the
central precision is $aI$, and the high precision is $P_rI$, exactly
as in \eqref{inv-zero-mode}. No scalar constant Laplacian was inverted.

For $M=1$, $Q=I$, $R=0$, and
\[
 \mathcal D=1+P_0/a,\quad k=\mathcal E_a=a/(P_0+a),\quad
 F=(P_0+a)^{-1}I,\quad\beta=d^+/(P_0+a).
\]
Thus
\begin{equation}\label{inv-M1}
 G_0=(P_0+a)^{-1}I+\frac{d^+d^-}{a(P_0+a)},
\end{equation}
the direct inverse of $(P_0+a)I-d^+d^-$. This checks the sign,
penalty normalization, and rank-one coefficient. The value $a^{-1}I$
at zero also shows why a uniform lower bound for $a$ is essential.

\subsection{Alias bounds with two fine derivatives}

On the common strip,
\[
 |P_r|\ge c|r|^2\ (r\ne0),\quad
 |u_{r,\mu}^{\pm}|\le C(1+|r|),\quad
 \sum_r|a_r^\pm|^2+
       \sum_{r,\mu}|\omega_{r,\mu}^\pm|^2\le C_d.
\]
The high-high block of $F$ is a bounded diagonal operator plus one
rank-one term for each orientation. Its high rank vectors are
$\omega_{r,\mu}^{\pm}/P_r$, and its mixed row and column use the
same vectors. Their alias $L^2$ norms are bounded. The central block
is bounded separately. Moreover \eqref{inv-beta-high} gives
\begin{equation}\label{inv-beta-estimate}
 |\beta_{r,\mu}|\le
 C\left(\frac{|a_r^-|}{|r|^3}
             +\frac{|\omega_{r,\mu}^-|}{|r|^2}\right),\qquad r\ne0.
\end{equation}
The central beta is bounded by \eqref{inv-beta-central}.
Therefore $F$ and $G_0$ have uniformly bounded alias Hilbert operator
norms. The norm estimate for an analytic rank-one matrix $v w^T$
uses $\norm{v}_2\norm{w}_2$, without conjugating its analytic parameters.

With $\alpha$ differences on the output and $\beta$ on the input,
$\alpha+\beta\le2$, the high diagonal of $F$ grows at most as
$|r|^{\alpha+\beta-2}$, which is bounded. Each rank vector has at
most two extra powers of $|r|$ multiplying
$\omega_{r,\mu}^{\pm}/P_r$, so its $L^2$ norm is still bounded by
Parseval. The mixed blocks have the same property. For the update in
\eqref{inv-final-formula}, two derivatives on a beta column are bounded,
in the worst case, by $C(|a_r^-|/|r|+|\omega_{r,\mu}^-|)$.
This is square summable uniformly, as is the corresponding row.
Every derivative sandwich of total order at most two is consequently
holomorphic and uniformly bounded in alias operator norm. The number
of derivative-orientation components is fixed in $d$.

\subsection{Reconstruction, periodic gluing, and the actual inverse}

For within-cell sites $s\in I_M^d$ use
$\mathcal F_M(z)_{sr}=e^{i(z+2\pi r)\cdot s/M}$.
For real $z$ it is unitary from alias counting $L^2$ into the cell
norm $M^{-d}\sum_s|v_s|^2$. For complex $z=x+iy$, it differs from
that unitary by $\diag_s(e^{-y\cdot s/M})$. Its norm and the norm of
its holomorphic inverse are thus uniformly bounded. Tensoring with
the orientation identity, the reconstructed symbol is
\[
 (G_0)_{\rm cell}(z)=\mathcal F_M(z)G_0(z)\mathcal F_M(z)^{-1}.
\]
All derivative sandwiches inherit its uniform strip bounds.

The full real inverse, rather than each separated central term, is
periodic up to alias permutation. Reconstruction cancels the
permutation; wrapping an alias gives $e^{i2\pi s_\nu}=1$. The
reconstructed symbol is therefore periodic. The enlarged-chart
holomorphic formulas agree on real overlaps and glue by the identity
theorem. Fourier contour shifting proves \eqref{inv-kernel-bound}.
Summing at a smaller exponent proves the weighted Schur bounds.
These are operator-valued contour estimates, so no factor growing
with the within-cell dimension is introduced.

For each fixed $M$, fiberwise multiplication, with the zero sector
as in \eqref{inv-zero-mode}, verifies that this bounded kernel is
the inverse of \eqref{inv-fine-precision} on the infinite homogeneous
Hilbert space. The real inverse bound yields a uniform positive fine
spectral gap. Sampling the same symbol on a compatible periodic
momentum grid gives exactly the finite source inverse. Absolute
summability identifies its kernel with the periodization; periodic
distance bounds its weighted sums by the infinite sums. This includes
short periods with their proper oriented-cell multiplicities.

Finally, at fixed coarse datum $B$, the Maxwell energy is at least
$\langle B,q_MB\rangle$, the source gauge term is nonnegative, and
the penalty is the constant $a\norm{B}^2$. The source-gauge minimizer
$H_{\rm src}B$ attains both minima. Therefore
\[
 \inf_{QA=B}\langle A,\mathcal A_aA\rangle
   =\langle B,(q_M+aI)B\rangle.
\]
The constrained Gaussian inverse formula proves
$q_0=q_M+aI$ and, by strict positivity and uniqueness,
$G_0Q^*q_0=H_{\rm src}$. In the source convention
$QH_0=\ell^{-1}$, the normalized extension is $H_0\ell$.
This proves \eqref{inv-coarse-result} and completes the theorem.

\begin{remark}\label{inv-scope}
The homogeneous source passage fixes $a=\ell^{-2}$ and hence $a=1$
here. Raw I, lines 23879--23881, separately introduces an $a_n$
sequence and refers to CMP99b (3.23)--(3.24) for an unprinted recursion.
The theorem does not guess that sequence: an application using it
must verify the fixed positive interval after its measure and scale
conversion. The scalar penalty $a'$ in the projection remains a
different object. General nonflat and retained-history identifications
are not consequences of the homogeneous Fourier proof.
\end{remark}
\section{The scalar gauge inverse and bounded current vertices}
\label{gauge-section}

The fine inverse alone does not bound the current perturbation.
The source gauge contains derivatives, so we first control its
scalar constrained inverse in the same physical cell norms.
Keep the matched scalar average $Q_0$, scalar Laplacian
$\mathcal L=D^*D$, and the source operators
\[
 G'=(\mathcal L+Q_0^*a'Q_0)^{-1},\qquad
 R=I-G'Q_0^*(Q_0G'^2Q_0^*)^{-1}Q_0G',
 \qquad Z=G'RG'.
\]
All adjoints use the physical measures. Complex continuations use
the fixed bilinear transpose, never conjugation of a varying parameter.

\subsection{The constrained biharmonic inverse}

The operator
\[
 Z=G'^2-G'^2Q_0^*(Q_0G'^2Q_0^*)^{-1}Q_0G'^2
\]
is the inverse of the quadratic form $\norm{\mathcal L\lambda}^2$
on $\ker Q_0$, extended by zero on its orthogonal complement.
Indeed it is the usual constrained inverse of $G'^{-2}$, and for
$Q_0\lambda=0$ one has
$\norm{G'^{-1}\lambda}^2=\norm{\mathcal L\lambda}^2$.
It follows that $Q_0Z=0=ZQ_0^*$ and that $Z$ is independent of
the auxiliary positive penalty $a'$. The constant scalar mode is
excluded by the constraint, so no inverse is assigned to that mode.

On a real nonzero Bloch fiber the constrained inverse is
\[
 Z_{rs}=\frac{\delta_{rs}}{P_r^2}
 -\frac{a_r^-a_s}{P_r^2P_s^2W_2},
 \qquad W_2=\frac{D_{s,2}}{P_0^2}.
\]
Its completely canceled form is
\begin{align}
 Z_{00}&=\frac{R_{s,2}}{D_{s,2}},
 &Z_{0s}&=-\frac{a_0^-a_s}{P_s^2D_{s,2}}\quad(s\ne0),
 \label{gauge-Z-central}\\
 Z_{r0}&=-\frac{a_r^-a_0}{P_r^2D_{s,2}}\quad(r\ne0),
 &Z_{rs}&=\frac{\delta_{rs}}{P_r^2}
       -\frac{P_0^2a_r^-a_s}{P_r^2P_s^2D_{s,2}}
          \quad(r,s\ne0).
 \label{gauge-Z-high}
\end{align}
The notation $D_{s,2}$ denotes the scalar denominator, not a
dependence on the second matrix index $s$. No central inverse
power of $P_0$ remains. At zero coarse momentum the high averaging
coefficients vanish: $Z$ is zero on the constant mode and equals
$P_r^{-2}$ on the nonzero aliases. For $M=1$, it is identically zero.

\begin{proposition}\label{gauge-Z-locality}
Every placement of at most four total fine forward or backward
differences on the two sides of $Z$ has a uniformly bounded
holomorphic physical-cell symbol and a uniformly exponentially
decaying cell-operator kernel. The bounds are independent of
refinement and compatible periodic volume.
\end{proposition}

\begin{proof}
Let the output order be $j$ and the input order be $l$, with
$j+l\le4$. The high diagonal in \eqref{gauge-Z-high} is bounded
by $C|r|^{j+l-4}$. Its high rank vectors are bounded by
$C|a_r^-||r|^{j-4}$ and $C|a_r||r|^{l-4}$.
Each order is at most four, so alias Parseval bounds both vector
Hilbert norms uniformly. The central row and column use the same
vectors; central multipliers are bounded. The scalar denominator
$D_{s,2}$ stays uniformly away from zero on the common strip.
For a complex rank-one matrix $vw^T$, its Hilbert norm is bounded
by $\norm v_2\norm w_2$, without conjugating the analytic symbol.

These are operator bounds before reconstruction and hence do not
count the growing number of aliases. The finite Fourier reconstruction
and its inverse are uniformly bounded in physical cell Hilbert
norm, with the inverse containing its required $M^{-d}$ factor.
Alias permutation and the integer-site phases give periodic gluing
on the enlarged charts. Contour shifting gives exponential cell
decay; summing at a smaller exponent gives weighted Schur bounds.
Fourier sampling gives exact periodizations on finite tori, including
the stated zero mode and small-period incidence multiplicities.
\end{proof}

The source definitions give
\[
 Q_0G'R=0,\quad \mathcal L G'R=R,\quad
 RG'\mathcal L=R.
\]
Consequently
\[
 \Lambda:=G'RD^*=Z\mathcal L D^*,\qquad
 D\Lambda=DZ\mathcal L D^*.
\]
The total fine derivative orders are three and four, respectively.
Proposition~\ref{gauge-Z-locality} therefore bounds both operators
in a common positive weighted cell-Schur norm. This is a Hilbert
operator conclusion; it does not assert pointwise bounds for four
fine derivatives.


\section{A uniform neighborhood for the flat auxiliary current}
\label{current-section}

We keep the background equal to the identity and vary the current
independently. Fine spacing is $M^{-1}$, with odd $M$, and the fine
measure is $M^{-d}$. The one-form average is the corrected source
average at flatness, already identified with $Q=Q_M^{\mathrm{st}}$.
The scalar source projection is $R$, and $D$ is the fine scalar
gradient. Write
\[
 \Delta_0=d_\eta^*d_\eta+DRD^*+Q^*aQ,
 \qquad G_0=\Delta_0^{-1},
 \qquad q_0=(QG_0Q^*)^{-1}.
\]
The homogeneous source convention has $a=1$ at unit physical block
scale. All conclusions below also hold uniformly for a fixed compact
interval $0<a_-\le a\le a_+<\infty$. We do not assert a bound for an
unmatched sequence of penalties.

The preceding flat inverse theorem and minimizing-extension theorem give
\begin{equation}
 q_0=q_M+aI,
 \qquad E_0=G_0Q^*q_0,
 \qquad QE_0=I,
 \qquad RD^*E_0=0.
 \label{current-base-identities}
\end{equation}
Here $q_M$ is the unpenalized constrained Maxwell precision. The
addition of $aI$ must be retained in $q_0$. The normalized extension
$E_0$ is unchanged by this penalty; a source convention
$QH_0=\ell^{-1}$ instead uses $H_0=E_0\ell^{-1}$.

All adjoints initially use these physical measures at the real flat
reference. At complex current they denote the fixed bilinear
transpose. They do not conjugate the varying complex current.

\subsection{Current vertices and the weighted algebra}
\label{current-vertices-section}

For a fine bond $b$, the flat gauge expansion has mixed and polarized
quadratic terms
\begin{align}
 c^{(1)}(A,\lambda)(b)
 &=\frac{i}{2}[\lambda(b_-)+\lambda(b_+),A(b)],
 \label{current-c1}\\
 2c^{(2)}(\lambda_1,\lambda_2)(b)
 &=\frac{1}{2i}\bigl(
 [\lambda_1(b_-),(D\lambda_2)(b)]
 +[\lambda_2(b_-),(D\lambda_1)(b)]\bigr).
 \label{current-c2}
\end{align}
These are the deposited formulas at Raw I:13698--13718 and
I:13811--13821. Define the current-linear maps by
\[
 \langle\lambda,\mathfrak j_J A\rangle
   =\langle c^{(1)}(A,\lambda),J\rangle,
 \qquad \mathfrak k_J=\mathfrak j_J^*,
\]
\[
 \langle\lambda_1,\mathfrak m_J\lambda_2\rangle
   =2\langle c^{(2)}(\lambda_1,\lambda_2),J\rangle.
\]
All bond and site pairings carry the same fine volume weight.
Set $\Lambda=G'RD^*$, with the fixed scalar source inverse $G'$.
The scalar locality lemma gives uniform weighted bounds for
$\Lambda$ and $D\Lambda$. The auxiliary perturbation is
\begin{equation}
 \Theta'_J=\mathfrak k_J\Lambda+
       \Lambda^*\mathfrak j_J+\Lambda^*\mathfrak m_J\Lambda.
 \label{current-theta-definition}
\end{equation}
This is precisely the auxiliary perturbation in Raw I:13725--13731.
It is distinct from the additional nonlinear-constraint operator
denoted there by $\Theta_J$.

Group fine fields into unit coarse cells with their physical Hilbert
norms. For block operators between these fine fibers, or between
fine and fixed-dimensional coarse fibers, use
\begin{equation}
 \norm{T}_\kappa=
 \max\left\{
   \sup_x\sum_y e^{\kappa d(x,y)}\norm{T_{xy}},
   \sup_y\sum_x e^{\kappa d(x,y)}\norm{T_{xy}}
 \right\}.
 \label{current-schur-norm}
\end{equation}
Here $d(x,y)=|x-y|_1$ in infinite volume and is periodic coarse-cell
distance on a finite compatible torus. This norm is submultiplicative,
including rectangular products, and transpose preserves it.

\begin{lemma}[A current bound without derivatives of the current]
\label{current-theta-bound}
There is $\theta<\infty$, independent of refinement and compatible
periodic volume, such that
\begin{equation}
 \norm{\Theta'_J}_\kappa\le\theta\norm{J}_\infty.
 \label{current-theta-norm}
\end{equation}
The map $J\mapsto\Theta'_J$ is complex linear. The current requires no
regularity beyond boundedness.
\end{lemma}

\begin{proof}
Define the finite-incidence map $V_J$ from fine bonds to fine sites by
\[
 \langle\lambda,V_J A\rangle
   =\frac{1}{2i}\langle[\lambda(b_-),A(b)],J(b)\rangle.
\]
Equation \eqref{current-c2} gives the exact form identity
\[
 \mathfrak m_J=V_JD+D^*V_J^*.
\]
Thus \eqref{current-theta-definition} becomes
\begin{equation}
 \Theta'_J=
 \mathfrak j_J^*\Lambda+\Lambda^*\mathfrak j_J
 +\Lambda^*V_JD\Lambda+(D\Lambda)^*V_J^*\Lambda.
 \label{current-theta-factorization}
\end{equation}
If $|[X,Y]|\le c_{\rm br}|X||Y|$, both
$\mathfrak j_J$ and $V_J$ have fine Hilbert operator norm at most
$C_d c_{\rm br}\norm{J}_\infty$. Indeed there are at most $2d$
incident bonds per site and two endpoints per bond; Cauchy--Schwarz
and the common physical volume weight give this bound without a
power of $M$. These coefficient maps have a fixed finite cell range,
so the same estimate holds in \eqref{current-schur-norm}.

Apply the weighted algebra inequality to
\eqref{current-theta-factorization}, using the established bounds for
$\Lambda,D\Lambda$. Every derivative is on a scalar leg through
$D\Lambda$. No derivative acts on $J$, even if $J$ jumps at a cell
boundary. The proof also applies to complex currents with the fixed
bilinear transpose.
\end{proof}

Choose a common sufficiently small $\kappa>0$. The flat inverse,
averaging and Maxwell locality results give constants $g,b,q$ with
\begin{equation}
 \norm{G_0}_\kappa\le g,
 \qquad \norm{Q}_\kappa,\norm{Q^*}_\kappa\le b,
 \qquad \norm{q_0}_\kappa\le q.
 \label{current-input-bounds}
\end{equation}
The fine inverse result also gives all physical fine-difference
sandwiches of total order at most two. These are established
homogeneous flat inputs; they are not hypotheses about a nonflat
source carrier. None of these constants includes a factor for the
number of sites in a fine cell.

\subsection{The scalar identity behind the zero-current response}

The source gauge of $E_0=H_{\rm src}$ gives $\Lambda E_0=0$.
For $B_0=G_0-E_0q_0^{-1}E_0^*$ one has
\begin{equation}\label{gauge-response-identities}
 G_0\Lambda^*=B_0\Lambda^*=DZ,\qquad \Lambda DZ=Z.
\end{equation}
Indeed, for $\lambda=Zf$ the scalar identities give
$Q_0\lambda=0$ and $\mathcal L\lambda=RG'f$, so
\[
 \Delta_0D\lambda=DRG'f=\Lambda^*f,\qquad QD\lambda=0.
\]
The first equality follows by the full inverse; the constrained
subtraction vanishes because $E_0^*\Lambda^*=0$. Also
$\Lambda DZ=G'R\mathcal LZ=G'RG'=Z$.
These exact identities will identify the first extension response
and the second coarse-precision response below.
\subsection{Uniform inverses and the exact quadratic remainder}
\label{current-neighborhood-section}

Define
\begin{equation}
 c=gb^2q,\qquad e=gbq,\qquad t_0=g(1+c),\qquad
 \varepsilon=\frac{1}{4g\theta(1+c)}.
 \label{current-radius}
\end{equation}
If $\theta=0$ the perturbation vanishes and the current restriction is
unnecessary. Otherwise use the displayed positive radius.

\begin{theorem}[Uniform flat auxiliary-current neighborhood]
\label{current-neighborhood-theorem}
For every complex bounded current with
$\norm{J}_\infty<\varepsilon$, the operators
\begin{equation}
 G_J=(\Delta_0-\Theta'_J)^{-1},\qquad
 q_J=(QG_JQ^*)^{-1},\qquad E_J=G_JQ^*q_J
 \label{current-family}
\end{equation}
exist and depend holomorphically on $J$ in the weighted algebra. They
are the deposited auxiliary family, with
$G_J=G_\pi^J$, $q_J=\mathsf q_J$ and $E_J=H^J\ell$. Uniformly in odd $M$, compatible volume, and
$a\in[a_-,a_+]$,
\begin{equation}
 \norm{G_J}_\kappa\le2g,\qquad
 \norm{q_J}_\kappa\le2q,\qquad
 QE_J=I,\qquad\norm{E_J}_\kappa\le4e.
 \label{current-uniform-bounds}
\end{equation}
Let
\begin{equation}
 B_0=G_0-E_0QG_0=G_0-E_0q_0^{-1}E_0^*.
 \label{current-B0}
\end{equation}
Then the exact formulas are
\begin{align}
 E_J-E_0
   &=(I-B_0\Theta'_J)^{-1}B_0\Theta'_JE_0,
 \label{current-extension-remainder}\\
 q_J-q_0
   &=-E_0^*\Theta'_J(I-B_0\Theta'_J)^{-1}
                  B_0\Theta'_JE_0.
 \label{current-quadratic-remainder}
\end{align}
In particular,
\begin{equation}
 \norm{E_J-E_0}_\kappa\le2t_0\theta e\norm{J}_\infty,
 \qquad
 \norm{q_J-q_0}_\kappa\le
       2e^2\theta^2t_0\norm{J}_\infty^2.
 \label{current-remainder-bounds}
\end{equation}
\end{theorem}

\begin{proof}
Abbreviate $\Theta=\Theta'_J$. The first Neumann series gives
\[
 G_J=(I-G_0\Theta)^{-1}G_0,
 \qquad
 \norm{G_J-G_0}_\kappa\le2g^2\theta\norm{J}_\infty,
\]
because $\norm{G_0\Theta}_\kappa<1/4$.
Writing $W_J=QG_JQ^*$, we obtain
\[
 \norm{q_0(W_J-W_0)}_\kappa
 \le2qb^2g^2\theta\norm{J}_\infty
 <\frac{c}{2(1+c)}<\frac12.
\]
The second Neumann series, with its specified order, is
\[
 q_J=[I+q_0(W_J-W_0)]^{-1}q_0.
\]
This proves \eqref{current-uniform-bounds}. Both series converge in
the same Banach algebra and therefore prove holomorphy on the entire
stated complex ball.

The constrained covariance has
\[
 QB_0=B_0Q^*=0,\qquad
 B_0\Delta_0=I-E_0Q,\qquad
 \norm{B_0}_\kappa\le t_0.
\]
Apply $B_0$ to
$(\Delta_0-\Theta)E_J=Q^*q_J$ and use $QE_J=I$. This gives
$E_J-E_0=B_0\Theta E_J$, hence
\eqref{current-extension-remainder}. No commutation of $B_0$ with
$\Theta$ is used. The chosen radius has
$\norm{B_0\Theta}_\kappa<1/4$.

Multiplying the stationary equation by $E_0^*$ gives
$q_J=q_0-E_0^*\Theta E_J$. The flat source gauge implies
$\Lambda E_0=0$, and \eqref{current-theta-definition} then gives
\begin{equation}
 \Theta'_JE_0=\Lambda^*\mathfrak j_JE_0,
 \qquad E_0^*\Theta'_JE_0=0.
 \label{current-linear-cancellation}
\end{equation}
Substitution of \eqref{current-extension-remainder} proves
\eqref{current-quadratic-remainder}. Finally
$\norm{E_0}_\kappa\le e$ and
$\norm{(I-B_0\Theta)^{-1}}_\kappa\le2$ give
\eqref{current-remainder-bounds}.
\end{proof}

\subsection{Factorial derivatives and localized current legs}
\label{current-derivatives-section}

For a current in this ball put
\begin{equation}
 B_J=G_J-E_JQG_J,
 \qquad t=2g(1+4c).
 \label{current-BJ}
\end{equation}
Then $B_J$ is the symmetric constrained covariance and
$\norm{B_J}_\kappa\le t$. All directional derivatives below are
derivatives in current space, not fine spatial differences.

\begin{proposition}[Exact derivative chains]
\label{current-derivative-chains}
For $n\ge1$ and bounded current directions $J_1,\ldots,J_n$,
\begin{equation}
 \begin{split}
 D^nq_J[J_1,\ldots,J_n]
 =-\sum_{\sigma\in S_n}
 E_J^*\Theta'_{J_{\sigma(1)}}B_J
 \Theta'_{J_{\sigma(2)}}\cdots B_J
 \Theta'_{J_{\sigma(n)}}E_J.
 \end{split}
 \label{current-q-derivative-chain}
\end{equation}
Each summand contains $n-1$ factors $B_J$; for $n=1$ it is
$-E_J^*\Theta'_{J_1}E_J$. The corresponding formula for $D^nE_J$
is the sum of the ordered products
$B_J\Theta'_{J_{\sigma(1)}}\cdots
B_J\Theta'_{J_{\sigma(n)}}E_J$. Consequently
\begin{align}
 \norm{D^nq_J[J_1,\ldots,J_n]}_\kappa
 &\le n!(4e)^2\theta^n t^{n-1}
                   \prod_i\norm{J_i}_\infty,
 \label{current-q-factorial}\\
 \norm{D^nE_J[J_1,\ldots,J_n]}_\kappa
 &\le n!(4e)(t\theta)^n\prod_i\norm{J_i}_\infty.
 \label{current-E-factorial}
\end{align}
\end{proposition}

\begin{proof}
At a fixed $J$, repeat the stationary-equation argument for a
sufficiently small increment $K$. It gives the exact identities
\begin{align}
 B_{J+K}&=(I-B_J\Theta'_K)^{-1}B_J,
 \label{current-B-increment}\\
 E_{J+K}&=(I-B_J\Theta'_K)^{-1}E_J,
 \label{current-E-increment}\\
 q_{J+K}&=q_J-E_J^*\Theta'_K
                    (I-B_J\Theta'_K)^{-1}E_J.
 \label{current-q-increment}
\end{align}
The first is the resolvent identity for the quadratic form
restricted to $\ker Q$, with its inverse extended by zero on
$\ran Q^*$. Expanding the series and taking the coefficient
multilinear in the directions produces exactly the permutations in
\eqref{current-q-derivative-chain} and the extension formula. The
norm bounds follow from submultiplicativity. The argument is valid
at each point in the ball by using a sufficiently small increment;
it does not enlarge the current domain.
\end{proof}

At zero current $Dq_0[J_1]=0$ by
\eqref{current-linear-cancellation}. The earlier scalar gauge
identity, with $Z=G'RG'$, gives the more explicit specializations
\[
 DE_0[J_1]=DZ\mathfrak j_{J_1}E_0,
\]
\begin{equation}
 \begin{split}
 D^2q_0[J_1,J_2]
 ={}&-(\mathfrak j_{J_1}E_0)^*Z(\mathfrak j_{J_2}E_0)\\
    &-(\mathfrak j_{J_2}E_0)^*Z(\mathfrak j_{J_1}E_0).
 \end{split}
 \label{current-zero-Hessian}
\end{equation}
The first extension response is a fine gradient. It need not
preserve the original source gauge at nonzero current.

\begin{proposition}[Paths through the current supports]
\label{current-localized-paths}
Let $J_{x_i}$ be supported on fine bonds based in cell $x_i$.
There are constants $C_0,C_1$, independent of refinement, compatible
volume, source locations and $n$, such that
\begin{equation}
 \begin{split}
 &\norm{D^nq_J[J_{x_1},\ldots,J_{x_n}](b,c)}\\
 &\quad\le C_0 C_1^n\prod_i\norm{J_{x_i}}_\infty
 \sum_{\sigma\in S_n}
 \exp\!\left[-\kappa\left(
 d(b,x_{\sigma(1)})
 +\sum_{i=1}^{n-1}d(x_{\sigma(i)},x_{\sigma(i+1)})
 +d(x_{\sigma(n)},c)\right)\right].
 \end{split}
 \label{current-path-bound}
\end{equation}
Here $b,c$ are coarse cells with their fixed orientation/color
fibers. Repeated source cells are allowed. The extension
derivatives have the same path structure, with the appropriate
fine/coarse end fibers.
\end{proposition}

\begin{proof}
For a current supported in cell $x$, each coefficient vertex
$\mathfrak j_{J_x}$ or $V_{J_x}$ reads and writes only a fixed
neighboring-cell set of $x$, and its restricted operator norm is
$O(\norm{J_x}_\infty)$. Expand each perturbation in a chain using
\eqref{current-theta-factorization}. Between consecutive local
vertices are products of the form $A B_J C^*$, with the appropriate
scalar/bond domains and with $A,C$ among
$I,\Lambda,D\Lambda$. Such products have a uniform weighted Schur
bound and hence their $(x,y)$ block is bounded by a constant times
$e^{-\kappa d(x,y)}$. The end products use $A E_J$ and obey the
same statement.

There are at most four choices per perturbation and a fixed number
of neighboring cells per coefficient vertex. Their sizes and
bounded offsets are absorbed into $C_1^n$. Multiplying the block
bounds and then summing the permutations proves
\eqref{current-path-bound}. This uses physical block operator norms
before multiplying; it never counts the fine sites in a block.
It also avoids convolution of pointwise bounds at a borderline
exponential rate. Periodic distance obeys the same triangle
inequality, so the proof is unchanged on compatible tori.
\end{proof}

\subsection{Fine derivatives, finite volume, and the scope of the result}
\label{current-scope-section}

The proved full flat inverse theorem controls every specified fine
difference sandwich of total order at most two. Write
$\mathcal D_\alpha$ and $\mathcal D_\beta$ for fixed words of those
orders, with the appropriate orientation domains. The exact identity
\begin{equation}
 \begin{split}
 \mathcal D_\alpha G_J\mathcal D_\beta
 ={}&\mathcal D_\alpha G_0\mathcal D_\beta\\
 &+(\mathcal D_\alpha G_0)\Theta'_J
       (I-G_0\Theta'_J)^{-1}G_0\mathcal D_\beta
 \end{split}
 \label{current-spatial-derivative-transfer}
\end{equation}
therefore gives the same uniform weighted bounds when
$\alpha+\beta\le2$. Derivatives act on the outside propagator
factors, not on the current. The one-sided estimates also give
the corresponding output fine derivatives of $E_J$. These are
physical cell $L^2$ operator estimates; no pointwise or H\"older
carrier estimate is inferred from them.

The reference kernels on finite compatible tori are exact
periodizations of the homogeneous kernels. All subsequent arguments
take place in the weighted block algebra on that torus. They apply
to arbitrary bounded periodic currents, without requiring them to
be constant from cell to cell, and retain constants uniform even
for short periods. The infinite-volume formulas follow in the
same algebra. No passage through an unweighted growing-dimensional
matrix norm is used.

The result is a uniform complex neighborhood for the displayed
\emph{auxiliary} flat family, with a quadratic precision remainder,
factorial current derivatives, and paths through their current
supports. The full source inverse instead has
\[
 G_1^J=[\Delta_0-\Theta'_J-\Pi^*\Theta_J\Pi]^{-1},
 \qquad \Pi=I-D\Lambda.
\]
The separate nonlinear-constraint term $\Theta_J$ is not covered
here. Its first coarse-precision derivative can already contain
$-E_0^*\Theta^{[1]}_J E_0$. Thus the cancellation
\eqref{current-linear-cancellation} is not a statement that the
full physical precision has zero current derivative. Nonflat
backgrounds, retained multiscale boundaries, source pointwise
carrier norms, and the interacting ultraviolet limit remain
outside this theorem.
\section{What has been discharged and what is next}

The flat auxiliary-current family now has a verified weighted
complex domain, rather than only formal derivatives at zero.
The inverse and extension bounds can supply propagator legs for
subsequent estimates, and the ordered source paths identify where
current derivatives are spatially supported. This discharges a
specific homogeneous input to the response program in Notebook 03.

Notebook 11 now bounds the additional operator $\Pi^*\Theta_J\Pi$
in this weighted Hilbert norm along $M=L^j$. Its proof retains the
corrected nonlinear one-step map and composes its quadratic jets
with their physical scale factors; a different rooted one-step
average is not substituted. Nonflat coefficients also change the covariant derivatives,
scalar gauge, and averaging map. These changes require a matched
uniform domain, not just a finite-regulator analytic inverse.
Pointwise source carrier estimates and retained-history boundary
estimates remain separate. Neither the free ultraviolet limit of
Notebook 07 nor the current neighborhood here constructs an
interacting four-dimensional Yang--Mills measure.

\begin{thebibliography}{9}
\bibitem{repo}
\href{https://github.com/ethanwsussman/Yang--Mills-LLM-experiment/tree/e392ef45013a581664bab76b90867170b30fe10e}
{Pinned Yang--Mills-LLM-experiment account}, commit
\texttt{e392ef45013a581664bab76b90867170b30fe10e}.
The source passages used here are Raw I:13583--13593
(coarse penalty), 13638--13655 (extension normalization),
13698--13731 and 13811--13821 (current vertices and auxiliary inverse),
13774 (gauge warning), 23909--23944 (scalar average and projection),
and 24349--24373 (corrected one-form average).
\bibitem{balaban109} T. Ba\l aban,
\emph{Renormalization group approach to lattice gauge field theories I},
Commun. Math. Phys. 109 (1987), 249--301.
\href{https://scispace.com/pdf/renormalization-group-approach-to-lattice-gauge-field-3ygai0vxt1.pdf}
{Copy of the original article}.
\end{thebibliography}
\end{document}
